% ID: SEL-D0012-BB-E04
% ESTADO: cotejada. Segunda pasada visual completa y borrador compilado revisados.
% Fuente: Selectividad/matematicas_aplicadas/EXTRA_RESERVA_B_Modelo-6-Enunciados (Mat-Ap-2021).pdf
% SHA-256: 1752ec2f46b4ef2cc5d1c3346d69b79c2cc8a7aa12f6a0cd8f0d994b5bcda157
% Página: 1; bloque B; ejercicio original 4.

La cotización en bolsa de una empresa en un determinado día viene expresada,
en euros, por la función $c(t)$, con $t\in[0,24]$, medido en horas.

La variación instantánea de esta función es la derivada de $c$, que viene dada
por $c'(t)=0.03t^2-0.9t+6$, con $t\in(0,24)$.
\begin{enumerate}[label=\alph*)]
  \item \textbf{(1.25 puntos)} Estudie los intervalos en los que la función
  $c$ es creciente.
  \item \textbf{(0.5 puntos)} Analice los puntos críticos de la función $c$,
  indicando en cuáles se alcanza el máximo y el mínimo relativos.
  \item \textbf{(0.75 puntos)} Halle la expresión analítica de la función
  $c$, sabiendo que la cotización en bolsa de la empresa era de 50 euros
  en el instante inicial.
\end{enumerate}
